
\lesson{23/11}
\section{Lorentz Group}
Recall some properties and definitions of the Lorentz group
(the defining relation using $\mathbf{(+,-,-,-)}$ metric, proper vs improper,
orthochronous vs non-orthochronous, time-like vs space-like vs light-like 
intervals).

\textbf{Examples of matrices of the Lorentz group}\\
\begin{enumerate}
\item Proper spatial rotations.
    Matrices in the form
        \[
            \begin{pmatrix}
                1&0&0&0\\
                0&&&\\
                0&&R&\\
                0&&&\\
            \end{pmatrix}\;\text{with }R\in SO(3)
        \]
        satisfy the definition of the Lorentz group. In particular, consider
        rotations about the $z$ axis:
        \[
            \begin{pmatrix}
                1&0&0&0\\
                0&\cos\vartheta&\sin\vartheta&0\\
                0&-\sin\vartheta&\cos\vartheta&0\\
                0&0&0&1\\
            \end{pmatrix}\;\text{with }\vartheta\in\left[0,2\pi\right]:
        \]
        for small $\vartheta$, this is
        \[
            \begin{pmatrix}
                1&0&0&0\\
                0&1&\vartheta&0\\
                0&-\vartheta&1&0\\
                0&0&0&1\\
            \end{pmatrix}+ O(\vartheta^2)= 
            \mathbb{I}+i\vartheta\underbrace{\begin{pmatrix}
                0&0&0&0\\
                0&0&-i&0\\
                0&i&0&0\\
                0&0&0&0
        \end{pmatrix}}_{:=J^3}+O(\vartheta^2).
        \]
    $J^3$ is the generator of rotations about the $z$ axis.

    Similarly,
    \[
        J^1 = \begin{pmatrix}
            0&0&0&0\\
            0&0&0&0\\
            0&0&0&-i\\
            0&0&i&0
        \end{pmatrix}\;J^2=
        \begin{pmatrix}
            0&0&0&0\\
            0&0&0&i\\
            0&0&0&0\\
            0&-i&0&0
        \end{pmatrix} 
    \]
    generate, respectively, rotations about $x,y$ axis.

    These three generators \textbf{satisfy the} $SU(2)$ \textbf{algebra}:
    \[
        \left[J^a,J^b\right] = i\varepsilon^{abc}J^c
    \]
    \textbf{Note}: the $\Lambda$ made from proper rotations are connected to
    the identity and so can be generated from the algebra, while using improper
    rotations (i.e. $O(3)\setminus SO(3)$) produces $\Lambda$ that are not.
\item Boosts at velocity $v$: defining
\[
    \beta:=\frac{v}{c}\;\gamma = \frac{1}{\sqrt{1-\beta^2}}
\]
a boost along the $x$ axis is
\[
    \Lambda(v) = \begin{pmatrix}
        \gamma&-\beta\gamma&0&0\\
        -\beta\gamma&\gamma&0&0\\
        0&0&1&0\\
        0&0&0&1
    \end{pmatrix}.
\]
Since $-1<\beta<1$ and $\gamma\geq 1$, this kind of matrix can be rewritten in
terms of hyperbolic functions of a parameter $\eta$, called \textit{rapidity}:
\[
    \gamma=\cosh\eta,\;\beta\gamma=\sinh\eta \longrightarrow \eta =
    \text{arcsinh}(\beta\gamma)
\]
and becomes
\[
    \Lambda(\eta) = 
    \begin{pmatrix}
        \cosh\eta&-\sinh\eta&0&0\\
        -\sinh\eta&\cosh\eta&0&0\\
        0&0&1&0\\
        0&0&0&1
    \end{pmatrix}:
\]
for small $\eta$, this can be expanded in
\[
    \Lambda(\eta) = \begin{pmatrix}
        1&-\eta&0&0\\
        -\eta&1&0&0\\
        0&0&1&0\\
        0&0&0&1
    \end{pmatrix} + O(\eta^2) =\mathbb{I}+i\eta\underbrace{\begin{pmatrix}
        0&i&0&0\\
        i&0&0&0\\
        0&0&0&0\\
        0&0&0&0
    \end{pmatrix}}_{:=K^1} + O(\eta^2)
\]
and similarly for the other boosts, generated by
\[
    K^2 = \begin{pmatrix}
        0&0&i&0\\
        0&0&0&0\\
        i&0&0&0\\
        0&0&0&0
    \end{pmatrix},\; K^3\begin{pmatrix}
        0&0&0&i\\
        0&0&0&0\\
        0&0&0&0\\
        i&0&0&0
    \end{pmatrix}. 
\]
They \textit{do not} form a closed sub-algebra for the Lorentz algebra, since
\[
    \left[K^a,K^b\right] = -i\varepsilon^{abc}J^c.
\]
The mixed commutator is
\[
    \left[J^a,K^b\right] = i\varepsilon^{abc}K^c
\]
\end{enumerate}
The general infinitesimal form is
\[
    \Lambda\indices{^\mu_\nu} = g\indices{^\mu_\nu}+\omega\indices{^\mu_\nu}
\]
Imposing the Lorentz invariance of the metric
\[
    g_{\mu\nu} = g_{\alpha\beta}\Lambda\indices{^\alpha_\nu}\Lambda\indices{^\beta
    _\mu} = g_{\alpha\beta}\left(g\indices{^\alpha_\nu}+\omega\indices{^\alpha_\nu
    }\right)\left(g\indices{^\beta_\mu}+\omega\indices{^\beta_\mu
    }\right) = g_{\mu\nu} +\omega_{\mu\nu}+\omega_{\nu\mu} + O(\omega^2)\iff 
    \omega_{\mu\nu}=-\omega_{\nu\mu}
\]
which means \textit{the generators of the Lorentz group are antisymmetric}:
this makes sense, because $4\times4$ antisymmetric matrices have $6$ degrees of
freedom and a Lorentz transformation is uniquely determined by $3$ angles of
rotations and $3$ boost velocities (or $3$ rapidities).

There are also interesting discreet transformations:
\begin{itemize}
\item Time reversal:
    \[
        \Lambda=\begin{pmatrix}
            -1&0&0&0\\
            0&1&0&0\\
            0&0&1&0\\
            0&0&0&1
        \end{pmatrix}
    \]
    which is non-orthochronous, improper and idempotent (i.e. $\Lambda^2=\mathbb
    {I}$).
\item Parity
    \[
        \Lambda=\begin{pmatrix}
            1&0&0&0\\
            0&-1&0&0\\
            0&0&-1&0\\
            0&0&0&-1        
        \end{pmatrix}   
    \]
    which is orthochronous, improper and idempotent.
\end{itemize}

Since $J$'s and $K$'s mix up, it is useful to construct linear combinations of
them to rewrite the algebra: consider
\[
    N^a = \frac{J^a+iK^a}{2},\;M^a = \frac{J^a-iK^a}{2},
\]
it can be proven that \textbf{they both generate an SU(2) closed subalgebra}:
\[
    \begin{dcases*}
        \left[N^a,N^b\right] = i\varepsilon^{abc}N^c\\ 
        \left[M^a,M^b\right] = i\varepsilon^{abc}M^c\\
        \left[M^a,N^b\right] = 0
    \end{dcases*}.
\] 
These relations can be proven by using brute force substitution and the
algebras of the $J$'s and $K$'s.

Recall that for $SU(2)$ the quantum number $J\in\mathbb{Z}/2$ is so that the
representation has
$\dim = 2J+1$ and that the eigenvalue of the Casimir is $J(J+1)$: since
the algebra of the (proper) Lorentz group is given by two $SU(2)$ algebras,
its representations can be classified by $2$ quantum numbers $p,q$:
\begin{table}[H]
    \centering
    \begin{tabular}{|c|c|c|}\hline
        $p$&$q$&Particle name\\\hline
        $0$&$0$&Scalar\\\hline
        $1/2$&$0$&LH Weyl spinor\\\hline
        $0$&$1/2$&RH Weyl spinor\\\hline
        $1/2$&$1/2$&Dirac spinor\\\hline
        $1$&$0$&Self-dual, antisymm. tensor\\\hline
        $0$&$1$&Anti-self-dual, antisymm. tensor\\\hline
    \end{tabular}
\end{table}
In the context of QFT and elementary particle physics, fields (particles) need
to be Lorentz invariant and can be characterized in this way.

In QFT, as the name suggests, the fundamental mathematical objects are fields:
to work with fields, differential operators needs to be introduced;
given a field $\phi(x^\mu)$, one can define
\[
    L_{\mu\nu} = i (x_\mu\pd_\nu-x_\nu\pd_\mu),
\]
which include \textit{realizations of both rotation and boost generators}.

They satisfy the \textit{Lorentz Algebra}:
\[
    \left[L_{\alpha\beta},L_{\gamma\delta}\right] = i\left(
    g\indices{_\alpha_\delta}L\indices{_\beta_\gamma}
    - g\indices{_\alpha_\gamma}L\indices{_\beta_\delta}
    - g\indices{_\beta_\delta} L\indices{_\alpha_\gamma} +
    g\indices{_\beta_\gamma} L\indices{_\alpha_\delta}\right).
\]
The infinitesimal transformations of the coordinates are
\[
    \delta x^\nu = \frac{i}{2}\varepsilon^{\alpha\beta}L_{\alpha\beta}x^\nu
\]
and their corresponding finite transformations are just the \textit{exponential
map} of the infinitesimal ones:
\[
    U(\Lambda) = \exp\left[ \frac{i}{2}\varepsilon^{\alpha\beta}L_{\alpha\beta}
    \right]
\]
Scalar fields transform like
\[
    \phi(x^\mu)\to \phi'(x'^\mu)U(\Lambda)\phi(x^\mu)U^{-1}(\Lambda)
\]
while vector fields transform like
\[
    \phi^\alpha(x^\mu) \to \phi^\beta(x'^\mu) = \Lambda\indices{^\beta_\alpha}
    \left(U(\Lambda)\phi^\alpha(x^\mu)U^{-1}(\Lambda)\right)
\]
\subsection{Extension to the Lorentz-Poincaré group}
Physical laws should also be independent of the point of spacetime
they are formulated in: this translates to the fact that physical laws
should also be \textit{invariant under space-time translations}:
\[
    x'^\beta = \Lambda\indices{^\beta_\alpha}x^\alpha + a^\beta
\]
with $a^\beta$ being any constant 4-vector. This can be represented as
\[
    (\Lambda,a)
\]
This seems, trivially, the direct
product of the Lorentz group and the translations, but in reality it is
the \textit{semi-direct product} of these two groups:
\[
    (\Lambda_1,a_1)(\Lambda_2,a_2) = (\Lambda_1\Lambda_2,a_1+\Lambda_1a_2).
\]
The generators of translations are
\[
    P^\mu = i\pd^\mu
\]
and they commute
\[
    \left[P^\alpha,P^\beta\right]=0:
\]
adding this to the Lorentz algebra yields
\[
    \begin{dcases*}
        \left[L_{\alpha\beta},L_{\gamma\delta}\right] = i\left(
        g\indices{_\alpha_\delta}L\indices{_\beta_\gamma}
        - g\indices{_\alpha_\gamma}L\indices{_\beta_\delta}
        - g\indices{_\beta_\delta} L\indices{_\alpha_\gamma} +
        g\indices{_\beta_\gamma} L\indices{_\alpha_\delta}\right)\\
        \left[P^\alpha,P^{\beta}\right]=0\\
        \left[L_{\alpha\beta},P_{\gamma}\right] = i (g_{\alpha\gamma} P_\beta
        - g_{\beta\gamma} P_\alpha)
    \end{dcases*}.
\]
So the total number of generators is $4+6=10$.

The Casimir are
\[
    P^\alpha P_\alpha = g_{\alpha\beta}P^\alpha P^\beta = m^2\mathbb{I}
\]
where the eigenvalue $m^2$ can be thought to be the square of the 
\textit{mass of the particle}.
The other Casimir is built from the \textit{Pauli-Lubanski axial vector}
\[
    W^\alpha = \frac{1}{2}\varepsilon^{\alpha\beta\gamma\delta}P_\beta
    L_{\gamma\delta}
\]
which is invariant under parity; the Casimir is
\[
    W^\alpha W_\alpha = g_{\alpha\beta}W^\alpha W^\beta = -m^2J(J+1)\mathbb{I}
\]
with $J\in\mathbb{Z}/2$ is \textit{the spin of the particle}.

If $m=0$, then it turns out that
\[
    W^\alpha=hP^\alpha,\;\text{with }h=\pm 1 
\]
where $h$ is called \textit{helicity}.

\textbf{Note} Continuous non-compact groups \textit{do not} have
finite-dimensional, non-trivial, unitary and irreducible representations.

\textbf{Def.} Wigner defined an elementary particle as a quantum excitation of a
field that corresponds to a unitary irreducible representations of the
Lorentz-Poincaré group.
